\chapter{Fused and Unfused Integration Kernels}\label{ch:fused}

The per-particle operations of a step --- locate, interpolate, integrate, apply the boundary
condition (\cref{fig:pipeline}) --- are realized by two interchangeable implementations. This
chapter describes the two paths, which configurations take which, and the guarantee that relates
them. The distinction is an implementation matter and does not change the numerics of the preceding
chapters; we document it because it decides the cost of a step (\cref{ch:performance}).

\section{Two implementations of one step}\label{sec:fused-two}

In the \emph{fused} path the entire step for a particle --- the nearest-node search at every
evaluation site, the spatial and temporal interpolation, the integrator update and the containment
check --- is carried out within one just-in-time-compiled kernel that loops over the particles in
parallel, with no intermediate arrays (\cref{fig:fused}); two exceptions are the adaptive RK45 step,
which tests containment after its kernel, and the three-dimensional RK2 kernel, which receives its start
anchor from the batch search. In the \emph{unfused} path the same
operations are separate routines, called in turn from Python, with the intermediate quantities
passed between them as whole arrays: all stage-one positions, then all stage-one velocities, and so
on. Each of these passes is itself a parallel compiled loop --- the stage interpolations of RK2 and
RK4 feed the stage positions to the vectorised interpolators --- so the unfused path also uses every
core; it pays for the extra passes over memory and, on a time-dependent flow, for materialising the
velocity field at each stage time. The fused path exists for speed; the unfused path exists for
generality, covering every combination of options for which no fused kernel is provided. Without
Numba every step is unfused.

\begin{figure}[t]
\centering
\begin{tikzpicture}[node distance=3mm and 7mm]
  \node[flowbox,text width=26mm] (u1) {\footnotesize locate (all)};
  \node[flowbox,below=of u1,text width=26mm] (u2) {\footnotesize interpolate (all)};
  \node[flowbox,below=of u2,text width=26mm] (u3) {\footnotesize integrate (all)};
  \node[flowbox,below=of u3,text width=26mm] (u4) {\footnotesize boundary (all)};
  \draw[flowarrow] (u1)--(u2); \draw[flowarrow] (u2)--(u3); \draw[flowarrow] (u3)--(u4);
  \node[font=\footnotesize,above=1pt of u1] {\textbf{Unfused path}};
  \node[flowbox,fill=laggreen!10,draw=laggreen,text width=40mm,minimum height=20mm,anchor=east]
    (F) at ($(u2.west)!0.5!(u3.west)+(-30mm,0)$)
    {\textbf{Fused kernel}\\[3pt]\footnotesize per particle:\\\footnotesize locate $\cdot$ interpolate $\cdot$\\\footnotesize integrate $\cdot$ boundary\\[3pt]\footnotesize (one compiled loop)};
  \draw[{Stealth[]}-{Stealth[]},gray,thick] (F.east) -- node[above,font=\footnotesize,align=center]{bit-identical\\results} (F.east -| u2.west);
\end{tikzpicture}
\caption{The two implementations of a single step. The fused kernel (left) performs all
per-particle operations inside one compiled loop; the unfused path (right) runs the same operations
as separate parallel passes over all particles. Where both exist they produce bit-identical
trajectories.}
\label{fig:fused}
\end{figure}

\section{Which path applies}\label{sec:fused-which}

\Cref{tab:fused-kernels} lists the fused kernel families and the conditions under which the step
functions choose them; every other configuration runs on the unfused path. In short: the
two-dimensional least-squares runs are fused for every integrator and every temporal treatment, except
Euler with \conf{em}/\conf{m1} diffusion (whose combined stochastic step replaces it) and, with the exact
guard off, HEALPix pixel meshes that fail the density certificate;
two-dimensional barycentric runs on node-layout velocities are fused when the flow is time-dependent,
reconstructed linearly and without element tracking;
three-dimensional runs are fused for Euler with least-squares velocities and for RK2 with node-layout
barycentric velocities on a z-level mesh. Among the paths that stay unfused are three-dimensional
RK4, every three-dimensional sigma-mesh run, element-layout barycentric and edge least-squares
interpolation, steady barycentric runs, and the \conf{idw} and \conf{leastsq} vertical combines of the
three-dimensional RK2 kernel. The last one matters in practice: the default
\conf{w\_interp: leastsq} keeps a node-layout three-dimensional RK2 run unfused, and setting
\conf{w\_interp} to \conf{nearest} or \conf{barycentric} moves it onto the fused kernel.

\begin{table}[t]
\centering\footnotesize\setlength{\tabcolsep}{4pt}
\begin{tabular}{@{}>{\raggedright\arraybackslash}p{30mm}>{\raggedright\arraybackslash}p{21mm}>{\raggedright\arraybackslash}p{73mm}@{}}
\toprule
kernel family & integrators & conditions \\
\midrule
2-D least squares, steady & Euler, RK2, RK4 & \conf{uv\_interp: leastsq} (element layout) \\
2-D least squares, time-aware & Euler, RK2, RK4 & as above; RK2 and RK4 with the next snapshot, Euler with a velocity between two snapshots, blended per element \\
2-D least squares, PCHIP & RK2, RK4 & as above, with \conf{temporal\_reconstruction: pchip} \\
2-D barycentric, time-aware & RK2, RK4 & \conf{uv\_interp: barycentric}, node layout, next snapshot present, linear in time, no carried element \\
3-D Euler & Euler & z-level mesh, \conf{uv\_interp: leastsq}, \conf{w\_interp} \conf{nearest} or \conf{leastsq} \\
3-D RK2 barycentric & RK2 & z-level mesh, node layout, \conf{uv\_interp: barycentric}, \conf{w\_interp} \conf{nearest} or \conf{barycentric}, no carried element; steady, linear or PCHIP \\
adaptive & RK45 & two-dimensional, element-layout least squares; this scheme has no unfused twin \\
\bottomrule
\end{tabular}
\caption{The fused kernel families and the configurations that select them. ``Carried element''
refers to \conf{elem\_search: walk} (\cref{sec:loc-walk}): the barycentric kernels keep their own
search, so a walk run takes the unfused path, while the least-squares kernels run and ignore the
carried element.}
\label{tab:fused-kernels}
\end{table}

The three-dimensional RK2 kernel evaluates per particle: the stage-one anchor and layer, the
horizontal velocity by barycentric interpolation with the temporal blend applied at each vertex, and
the vertical velocity $W_1$ --- computed before the horizontal verdict, so that a particle that dies
on the horizontal interpolation reports the same vertical velocity on both paths; the midpoint
$\mathbf{x}_{1/2}$ and $z_{1/2}=z+W_1\,\Delta t/2$; the stage-two anchor, layer and velocities; the end
position with $z_{\text{new}}=z+W_2\,\Delta t$; the containment test; and the clamp of the depth to the
water column. At one million particles on a mesh of \num{196000} nodes and 70 levels it was measured
1.30 times faster than the unfused path at 16 threads (a development measurement not kept in the
repository); the gain grows with the mesh and field size, while
scattered per-particle reads dominate both paths (\cref{sec:perf-schemes}). The two-dimensional
barycentric kernels remove the per-stage materialisation of the whole node field: on a
four-million-node mesh an RK4 step fell from \SI{83}{ms} to \SI{9}{ms} at eight threads, independent of
the mesh size.

\section{Bit-identity}\label{sec:fused-identity}

Where a fused and an unfused implementation of the same configuration both exist, they produce
bit-identical trajectories, statuses and failure reasons, which the test suite checks
(\cref{ch:validation}). Three things make it hold.

\paragraph{The temporal blend.} Rather than materialising the field at each stage time, the fused
kernels apply a nested two-level blend inline,
\begin{equation}\label{eq:fused-nested}
  \mathbf{u}_k=(1-\alpha_k)\bigl[(1-\alpha_0)\,\mathbf{u}_A+\alpha_0\,\mathbf{u}_B\bigr]+\alpha_k\,\mathbf{u}_T ,
\end{equation}
where $A$ and $B$ are the snapshots bracketing the step's start, $\alpha_0$ is the start's position in
that bracket, and $T$ is the stage \emph{target} snapshot the step moves toward --- the later bracket
snapshot when integrating forward, the earlier one when integrating backward, which is why the target
is passed separately from the bracket. \Cref{eq:fused-nested} performs the same floating-point
operations as stepping on the materialised field of \cref{eq:temp-linear}; with $\alpha_0=0$ it
collapses to a single blend, and the barycentric kernels short-circuit the end-point fractions exactly
as the unfused stage construction does, so the identity holds for non-finite values too.

\paragraph{The anchors.} Every evaluation site of a fused kernel --- the start, each stage and the
final position --- finds its nearest node through the exact search of \cref{sec:loc-exact}, the same
function the unfused batch query calls, so the two paths anchor every particle to the same node, ties
included. This is the default (\conf{mesh.exact\_nearest: true}). With the guard off, the fused
kernels widen an empty window ring by ring --- an approximate anchor --- and identity rests on gates
instead: the two-dimensional barycentric kernels (and the least-squares kernels on HEALPix pixel
meshes) run only on meshes that pass the density certificate of \cref{sec:loc-exact}, with a warning
and a fall-back to the unfused path otherwise; the three-dimensional RK2 kernel flags a particle whose
stage window is empty, repairs the flagged positions with the batch search and re-runs only those
particles through the kernel with the repaired anchor, for up to two rounds, before a last-resort
unfused re-run of the remainder. The least-squares kernels on FESOM-family meshes and the fused
three-dimensional Euler kernel have no such gate: with the guard off they agree with the unfused path
only where no three-by-three window is empty, which holds in the domain interior of quasi-uniform meshes.

\paragraph{The diagnostics.} Each fused kernel records why a particle failed --- interpolation failure
or out of bounds --- in a per-particle scratch array at its failure branches, and the caller merges it
into the sticky failure-reason field (\conf{mode.track\_failure\_reason}, \cref{ch:lifecycle}) exactly
as the unfused path does. A failed particle keeps its start position and anchor on both paths, and
frozen particles are copied verbatim.

\codeanchor{The fused kernels are the \fn{\_fused\_*} functions of \srcfile{numerics/integrators.py};
their runners are \fn{\_run\_fused\_2d}, \fn{\_run\_fused\_euler\_2d\_pair},
\fn{\_run\_fused\_2d\_tinterp}, \fn{\_run\_fused\_2d\_bary\_tinterp}, \fn{\_run\_fused\_2d\_pchip},
\fn{\_run\_fused\_euler\_3d} and \fn{\_run\_fused\_rk2\_3d\_bary\_tinterp}/\fn{\_pchip}; the adaptive
scheme is \fn{\_rk45\_kernel\_2d}. The dispatch lives in \fn{euler\_step}, \fn{rk2\_step} and
\fn{rk4\_step}; the anchor of each site is \fn{\_fused\_site\_anchor}; the guard-off gates are
\fn{\_bary\_fused\_anchor\_exact} and \fn{\_leastsq\_fused\_anchor\_exact}, and the repair loop
\fn{\_run\_fused\_3d\_with\_heal}; failure reasons are merged by \fn{\_merge\_fused\_failure}. The
unfused stage interpolation goes through \fn{\_interp\_uv\_at}. The inline blends are the
\fn{\_interp\_*\_tinterp} and \fn{\_interp\_*\_pchip4} cores of \srcfile{numerics/interpolate.py}
and, for the vertical velocity, of \srcfile{numerics/vertical.py}.}
